\subsubsection{Tiêu chí đánh giá độ chính xác}
Độ chính xác của mô hình TinyML được đánh giá theo hai hướng bổ trợ cho nhau. Thứ nhất là đánh giá ngoại tuyến trong quá trình huấn luyện, sử dụng các chỉ số \textit{Accuracy} và \textit{Confusion Matrix} trên tập kiểm thử để xem mô hình phân biệt bốn lớp \texttt{Cloudy}, \texttt{Rainy}, \texttt{Snowy} và \texttt{Sunny} ở mức độ nào. Đây là bước đánh giá định lượng trong môi trường Python. Thứ hai là đánh giá trực tuyến trên thiết bị ESP32 thông qua kết quả suy luận hiển thị trên giao diện web và log chi tiết xuất ra terminal, nhằm kiểm tra xem mô hình sau khi chuyển sang TensorFlow Lite Micro có còn cho ra hành vi nhất quán hay không.

Cách tiếp cận này có ý nghĩa quan trọng vì một mô hình đạt kết quả tốt trên máy tính chỉ thực sự có giá trị khi vẫn giữ được tính nhất quán sau khi chuyển sang TensorFlow Lite và triển khai trên vi điều khiển. Do đó, ngoài việc quan tâm tới nhãn dự đoán cuối cùng, nhóm còn kiểm tra cả chuỗi xử lý trung gian gồm dữ liệu thô từ cảm biến, giá trị sau chuẩn hóa và phân bố xác suất ở đầu ra của mô hình. Nói cách khác, phần đánh giá trong báo cáo không chỉ dừng ở câu hỏi “mô hình có đúng không”, mà còn kiểm tra “mô hình có còn đúng sau khi được nhúng lên phần cứng hay không”.

\subsubsection{Nhận xét từ kết quả triển khai thực tế}
Qua các ảnh minh họa ở phần Task 5, có thể thấy mô hình TinyML đã vận hành đúng theo thiết kế. Trên giao diện web, hệ thống có khả năng hiển thị trực tiếp nhãn dự đoán thời tiết để người dùng quan sát một cách trực quan. Song song với đó, terminal cung cấp đầy đủ các thông tin kiểm chứng như nhiệt độ ban đầu, độ ẩm ban đầu, giá trị đã chuẩn hóa và xác suất của từng lớp. Cơ chế này giúp xác nhận rằng mô hình trên ESP32 không chỉ đưa ra một kết quả cuối cùng, mà còn thực hiện đầy đủ pipeline suy luận như ở môi trường huấn luyện.

Với đặc trưng đầu vào chỉ gồm hai biến là nhiệt độ và độ ẩm, mô hình vẫn cho thấy khả năng phân loại tốt trong các trường hợp đặc trưng rõ rệt. Chẳng hạn, những mẫu có độ ẩm cao thường nghiêng về lớp \texttt{Rainy}, trong khi các mẫu có nhiệt độ cao và độ ẩm thấp hơn thường có xu hướng thuộc lớp \texttt{Sunny}. Tuy nhiên, các lớp có vùng giá trị chồng lấn nhau như \texttt{Cloudy} và \texttt{Rainy} vẫn có thể tạo ra sai số nhất định. Đây là giới hạn tự nhiên của mô hình gọn nhẹ, đồng thời cũng là đánh đổi hợp lý để đổi lấy tốc độ suy luận nhanh và khả năng chạy trực tiếp trên thiết bị nhúng.

\subsubsection{Đánh giá tổng quan}
Nhìn chung, mô hình TinyML đáp ứng tốt mục tiêu của đề tài là đưa khả năng nhận biết trạng thái môi trường xuống ngay tại biên mạng. Điểm mạnh lớn nhất là tính gọn nhẹ, tốc độ phản hồi nhanh và khả năng tích hợp chặt chẽ với hệ thống giám sát hiện có. Dù chưa hướng đến một mô hình phân loại phức tạp với số lượng đặc trưng lớn, giải pháp hiện tại đã đủ để minh chứng tính khả thi của việc kết hợp AI với hệ thống IoT trên nền tảng ESP32.